\section*{写在前面}
\addcontentsline{toc}{section}{写在前面}

\subsection*{这份讲义是什么}

这是一份\term{观点讲义}：把某位博主在一系列视频里反复表达的世界观、概念与方法，提炼、重组、延伸成一份可以从头读到尾的连贯文档。它不是视频字幕的摘抄合集，而是一次再创作——素材属于博主，结构属于整理者。

\subsection*{这份讲义不是什么}

\begin{checklist}
  \item 不是学术综述，也不是专业建议。涉及医学、法律、投资等领域时，博主的经验只能作为视角，不能替代专业人士；
  \item 不是逐字稿。口语被改写成书面表达，例子被重新编排，但观点不替博主发明；
  \item 不是「正确答案」。凡是博主本人标注为猜测、个体经验的内容，讲义同样保留这份不确定。
\end{checklist}

\subsection*{怎么读}

\key{先接受一套新的解释框架，再决定拿走哪些方法。} 第一部分通常回答「为什么会这样」（认知层），中间部分回答「那怎么办」（方法层），最后是延伸思考与清单。可以顺序读，也可以直接跳到与你处境最相关的一章。
